\section{Conclusion}

We introduced CFR-WM, a compact world model that combines a prespecified,
mask-exact complementary context update with a probabilistic
physics-residual transition. The design isolates two sources of long-horizon
error: uncertainty in the state at the forecast origin and error accumulated by
the free-rollout transition. A conditional error identity explains when the
fixed update can help, while gain and training-horizon sensitivities test its
limits without claiming universal filtering optimality.

The confirmatory study trained \ConfirmatoryRuns{} models across five
controlled nonlinear systems, three deployment regimes, eight independent
seed-level realizations, five broad comparators, and five sensitivities. Against
the preselected PGNN-style baseline, the seed-balanced relative NRMSE-25 effect was
\PrimaryImprovement\% (95\% CI [\PrimaryRelativeCILow\%,
\PrimaryRelativeCIHigh\%]); accordingly, H1 was \PrimaryDecision, but the
prespecified H2 robustness claim was \RobustnessDecision{} because clean
noninferiority failed. Specifically, the corrupt/shifted component improved while clean-regime
noninferiority failed, exposing the cost of a nonadaptive equal-trust update.
These decisions follow the frozen 5\% practical margin and exact seed-level
tests, not a favorable pooled mean or selected system.

The work contributes both a technical composition and an evidential pattern:
simple physics-aware corrections should be evaluated against matched
sensitivities, and unsuccessful development mechanisms should remain visible
rather than being absorbed into confirmation. The present evidence is limited
to state-transition prediction with known controls on a synthetic-mechanistic
benchmark. The next defensible steps are preregistered evaluation with
structurally misspecified physics and informative missingness, followed by
visual state estimation and hardware-in-the-loop closed-loop control. Those
extensions require new safety, calibration, and policy endpoints; they cannot
be inferred from NRMSE alone.
